%

\section{Twierdzenie Cevy}

Ważnym kryterium współpękowości trzech czewian jest:

\begin{theorem}[Cevy]
\label{twierdzenie_cevy}%
	Dany jest trójkąt $ABC$ i trzy różne od wierzchołków punkty $K \in BC$, $L \in CA$, $M \in AB$.
	Wówczas czewiany $AK$, $BL$, $CM$ są współpękowe wtedy i tylko wtedy, gdy
	\begin{equation}
		[AMB] [BKC] [CLA] = 1.
	\end{equation}
	\index{twierdzenie!Cevy}
\end{theorem}
% Neugebauer 119
% Ceva = guzicki-3
% Eves s. 248

https://www.deltami.edu.pl/2021/09/zabawy-na-polu/ 7

Twierdzenie należy do geometrii afinicznej, ponieważ można je wyrazić bez odnoszenia się do kątów, pól albo długości.
Przypisze się je Giovanniemu Cevie, który wyda pracę \emph{De lineis rectis} (o~liniach prostych) w 1678 roku; chociaż twierdzenie nie będzie obce Yusufowi al-Mu'tamanowi ibn Hudzie, emirowi taify w Saragossie z XI wieku.
\index[persons]{al-Mutaman, Yusuf ibn Huda}%
\index[persons]{Ceva, Giovanni}%
Wspomniana praca Cevy będzie zawierać także twierdzenie Menelaosa, odkryte na nowo przez Cevę.

Piszą o nim (twierdzeniu, nie emirze) Zetel \cite[s. 11-14, 40]{zetel_2020}, Audin \cite[s. 38]{audin_2003}, Hartshorne \cite[s. 182]{hartshorne2000}.
Zetel stwierdzi, że uogólnienie twierdzenia Cevy nazywa się twierdzeniem Ponceleta, mimo braku podobieństw do faktu \ref{big_poncelet}.
Osiem zadań, do rozwiązania których wygodnym narzędziem jest twierdzenie Cevy, poda Joanna Jaszuńska w $\Delta_{11}^2$. % https://www.deltami.edu.pl/2011/02/twierdzenie-cevy/

Wnioskiem z twierdzenia Cevy jest wynik znany jako lemat Steinbarta:

\begin{lemma}[Steinbarta]
	W trójkąt $ABC$ wpisano okrąg, styczny do boków $BC$, $AC$ i $AB$ kolejno w punktach $D$, $E$, $F$.
	Punkty $X$, $Y$, $Z$ leżą odpowiednio na krótszych łukach $EF$, $FD$ i $DE$.
	Następujące warunki są równoważne:
	\begin{itemize}
		\item proste $XD$, $YE$, $ZF$ są współpękowe,
		\item proste $AX$, $BY$, $CZ$ są współpękowe.
	\end{itemize}
\end{lemma}

Piszą o nim Ambroszczyk, Komisarczyk \cite[s. 17]{komisarczyk_2024}, ale pełną historię wyniku znajdziemy w artykule
% https://www.tandfonline.com/doi/full/10.1080/0025570X.2025.2595954#d1e168

% @article{Humenberger01012026,
% author = {Hans Humenberger},
% title = {An Introduction to and a Further Variation of the Steinbart Theorem},
% journal = {Mathematics Magazine},
% volume = {99},
% number = {1},
% pages = {20--29},
% year = {2026},
% publisher = {Taylor \& Francis},
% doi = {10.1080/0025570X.2025.2595954},
% URL = { 
%         https://doi.org/10.1080/0025570X.2025.2595954
% },
% eprint = { 
%         https://doi.org/10.1080/0025570X.2025.2595954
% }
% }

% * The phenomenon described in the Steinbart theorem was already formulated 1991 as a problem by S. Rabinowitz ([4], of course, not calling it Steinbart theorem; only for points P in the interior of the triangle) and a solution was given by F. Bellot and M. A. Lopez, in the following year 1992, see [5]. If a reader happens to know an even earlier source of the phenomenon described in the Steinbart theorem, please inform the author.
% https://www.tandfonline.com/doi/full/10.1080/0025570X.2025.2595954#abstract
% The Steinbart theorem seems to be not very well known in the community, it was detected and proved by an ambitious German high school student, Oliver Funck.
% The high school student Oliver Funck (Steinbart gymnasium in Duisburg, Germany) detectedFootnote* and proved (with the help of a computer algebra system; below we present a geometric proof) this phenomenon in the year 2000 presenting a note (16 pp) in the framework of “Jugend forscht” (German for “youth researches,” see [Citation7]). To the honor of the school’s eponym, Quintin Steinbart, he called Q the Steinbart point of the triangle ABC with respect to point P and the corresponding theorem Steinbart’s theorem. This labeling should remind one of the Exeter point, a special case of the Steinbart point, namely 𝑃= centroid G of the triangle, detected in 1986 at a computers-in-mathematics workshop at Phillips Exeter Academy (this point lies on the Euler line). Since the Exeter point was named after the institution in which it was detected, Oliver Funck also named “his” point after “his” institution (Steinbart gymnasium).


Prostą, która przechodzi przez wierzchołek trójkąta i przeciwległy bok nazywa się czewianą, nazwy tej używa się też wobec odcinka trójkąta leżącego na tej prostej.

\begin{proposition}[wzór Steinera-Routha]
	Rozpatrzmy trójkąt $ABC$ z czewianami $AD$, $BE$, $CF$ tak, że
	\begin{equation}
		\frac{BD}{DC} = \lambda_1, \quad
		\frac{CE}{EA} = \lambda_2, \quad
		\frac{AF}{FB} = \lambda_3.
	\end{equation}
	Oznaczmy przekrój $BE$ i $CF$ przez $G_1$; $AD$ i $CF$ przez $G_2$, zaś $AD$ i $BE$ przez $G_3$.
	Wtedy
	\begin{equation}
		\frac{S_{G_1G_2G_3}}{S_{ABC}} = \frac{(\lambda_1 \lambda_2 \lambda_3 - 1)^2}{
			(\lambda_1 \lambda_2 + \lambda_1 + 1)
			(\lambda_2 \lambda_3 + \lambda_2 + 1)
			(\lambda_3 \lambda_1 + \lambda_3 + 1)
		}
	\end{equation}
	oraz
	\begin{equation}
		\frac{S_{DEF}}{S_{ABC}} = \frac{\lambda_1 \lambda_2 \lambda_3 + 1}{
			(\lambda_1 + 1)
			(\lambda_2 + 1)
			(\lambda_3 + 1)
		}.
	\end{equation}
\end{proposition}

% J. Steiner, Gesammelte Werke, Vol. 1, Reimer, Berlin, 1882. strony 163-168
% E. J. Routh, A Treatise on Analytical Statics with Numerous Examples, Vol. 1, second edition. Cambridge University Press, 1896. strona 82
% https://ems.press/content/serial-article-files/51635

Stanowi to uogólnienie twierdzeń Cevy ($\lambda_1 \lambda_2 \lambda_3 = 1$) i Menelausa ($\lambda_1 \lambda_2 \lambda_3 = -1$).
Napiszą o nim w Delcie ($\Delta_{11}^6$, $\Delta_{12}^3$, $\Delta_{16}^8$)

\begin{example}[trójkąt, który zaskoczył Feynmana]
	Dla $\lambda_1 = \lambda_2 = \lambda_3 = 2$, stosunek pól trójkątów wynosi $1 : 7$.
\end{example}

% https://www.deltami.edu.pl/2011/06/sprawa-niezbyt-pedagogiczna/
% https://www.deltami.edu.pl/2012/03/co-moga-nam-dac-ciezary-i-wypory/
% https://www.deltami.edu.pl/2016/08/polecamy-bardzo-trudne-zadanie/

Ceva 3D https://www.deltami.edu.pl/2013/10/kacik-przestrzenny-19-plaszczyzny-przecinajace-sie-w-punkcie/

\subsubsection{Punkt Nagela}
% eksperymentalne tłumaczenie na włoski

%

TODO: Punkt Nagela, pojawia się u Evesa \cite[s. 71]{eves1_1972} jako ćwiczenie na zastosowanie twierdzenia Cevy razem z punktem Gergonne'a. % lang-pl

\index{punkt!Nagela} % lang-pl
TODO: Il punto di Nagel compare in Eves \cite[p. 71]{eves1_1972} come esercizio sull'applicazione del teorema di Ceva insieme al punto di Gergonne. % lang-it

\index{czewiany Nagela} % lang-pl
\index{punto!di Nagel} % lang-it

\index{ceviane di Nagel} % lang-it

\begin{proposition}[punkt Nagela] % lang-pl
\begin{proposition}[punto di Nagel] % lang-it
    \label{punkt_nagela}
\end{proposition}

O punkcie Nagela pisze też Zetel \cite[s. 22, 25, 64]{zetel_2020} (punkt Nagela, punkt przecięcia środkowych -- środek ciężkości trójkąta, środek okręgu wpisanego oraz środek ciężkości obwodu trójkąta leżą na jednej prostej zwanej drugą prostą Eulera albo prostą Eulera-Nagela). % lang-pl

Del punto di Nagel scrive anche Zetel \cite[p. 22, 25, 64]{zetel_2020} (punto di Nagel, punto di concorrenza delle mediane: il baricentro del triangolo, l'incentro e il baricentro del perimetro del triangolo giacciono su una stessa retta, detta seconda retta di Eulero oppure retta di Eulero--Nagel). % lang-it

% % Punkty izogonalnie sprzężone w trójkącie % https://www.deltami.edu.pl/2010/11/izogonalnie-sprzezone/

\subsubsection{Punkt Gergonne'a}
% eksperymentalne tłumaczenie na włoski

%

TODO: Czewiany Gergonne'a. % lang-pl
\index{czewiany Gergonne'a} % lang-pl
TODO: Ceviane di Gergonne. % lang-it
\index{ceviane di Gergonne} % lang-it

\begin{definition}
[punkt Gergonne'a] % lang-pl
[punto di Gergonne] % lang-it
\index{punkt!Gergonne'a}% % lang-pl
\index{punto!di Gergonne}% % lang-it
\label{punkt_gergonne}
    Dany jest okrąg wpisany w~trójkąt $ABC$, styczny do boków $BC$, $AC$, $AB$ odpowiednio w~punktach $P$, $Q$ i $R$. % lang-pl
    Wtedy czewiany $AP$, $BQ$ i $CR$ są współpękowe. % lang-pl
    Sia dato il cerchio inscritto nel triangolo $ABC$, tangente ai lati $BC$, $AC$, $AB$ rispettivamente nei punti $P$, $Q$ e $R$. % lang-it
    Allora le ceviane $AP$, $BQ$ e $CR$ sono concorrenti. % lang-it
\end{definition}

TODO: rysunek % lang-pl
TODO: figura % lang-it

To samo, ale bez użycia słowa ,,czewiany'' napisze Zetel \cite[s. 15, 25]{zetel_2020}. % lang-pl
Lo stesso risultato, ma senza usare la parola ,,ceviane'', è presentato da Zetel \cite[p. 15, 25]{zetel_2020}. % lang-it

%

\subsubsection{Punkt Lemoine'a}
% eksperymentalne tłumaczenie na włoski

%

\index{punkt!Lemoine'a} % lang-pl
\index{punto!di Lemoine} % lang-it

Emile Lemoine odczyta w 1873 roku pracę \emph{,,Sur quelques propriétés d'un point remarquable du triangle''}, gdzie po raz pierwszy pojawią się bez dowodów fakty \ref{random_prp_symmedians_proportional_squares} i \ref{punkt_lemoine}. % lang-pl
Termin symediana zaproponuje Maurice d'Ocagne w 1883 roku, jako zastępstwo \emph{antyrównoległej środkowej} używanego przez Lemoine'a. % lang-pl
Nel 1873 Émile Lemoine presenterà il lavoro \emph{,,Sur quelques propriétés d'un point remarquable du triangle''}, nel quale compariranno per la prima volta, senza dimostrazione, i fatti \ref{random_prp_symmedians_proportional_squares} e \ref{punkt_lemoine}. % lang-it
Il termine \emph{symmediana} sarà proposto da Maurice d'Ocagne nel 1883 come sostituto di \emph{mediana antiparallela}, espressione usata da Lemoine. % lang-it
% Nouvelles Annales de Mathématiques, 3rd series, II, 451 (1883).
Punkt Lemoine'a pojawia się u Józefa Neuberga (praca \emph{Mémoire sur le tétraèdre} z 1886 roku). % lang-pl
Il punto di Lemoine compare nell'opera di Joseph Neuberg (\emph{Mémoire sur le tétraèdre}, 1886). % lang-it
% https://www.persee.fr/doc/marb_0770-8459_1886_num_37_1_2317
W~1847 roku Ernst Grebe skonstruuje kwadraty na zewnątrz trójkąta i odkryje punkt Lemoine'a przed Lemoinem, dlatego w literaturze padać będzie też nazwa punkt Grebego. % lang-pl
Nel 1847 Ernst Grebe costruirà quadrati all'esterno di un triangolo e scoprirà il punto di Lemoine prima dello stesso Lemoine; per questo motivo nella letteratura compare anche il nome di punto di Grebe. % lang-it

Wszystko to pięknie opisze Altshiller-Court \cite{altshiller_2007}. % lang-pl
Tutto ciò è descritto magnificamente da Altshiller-Court \cite{altshiller_2007}. % lang-it

% UW: Twierdzenie Cevy (wraz z trygonometryczną wersją),
% przykłady punktów szczególnych trójkąta:
% punkt Nagela (Guzicki-4),
% punkt Gergonne'a (guzicki4),

% https://www.deltami.edu.pl/2016/05/izogonalne-sprzezenie-i-symediany/

\begin{definition}
[symediana] % lang-pl
[symmediana] % lang-it
    Odbicie symetryczne środkowej trójkąta względem dwusiecznej wychodzącej z~tego samego wierzchołka nazywamy symedianą. % lang-pl
    La riflessione della mediana di un triangolo rispetto alla bisettrice uscente dallo stesso vertice si chiama symmediana. % lang-it
\end{definition}

Symediana to prosta izogonalna ze środkową (Zetel \cite[s. 111]{zetel_2020}; Ambroszczyk i Komisarczyk \cite[s. 26]{komisarczyk_2024}; Dutka w $\Delta_{15}^2$). % lang-pl
La symmediana è la retta izogonale della mediana (Zetel \cite[p. 111]{zetel_2020}; Ambroszczyk e Komisarczyk \cite[p. 26]{komisarczyk_2024}; Dutka in $\Delta_{15}^2$). % lang-it
% https://www.deltami.edu.pl/2015/02/krotka-opowiesc-o-symedianie/
Symediana jest miejscem geometrycznym punktów, których odległości od dwóch boków trójkąta są proporcjonalne do tych boków. % lang-pl
Ale dużo ważniejszą jej własnością jest: % lang-pl
La symmediana è il luogo geometrico dei punti le cui distanze da due lati del triangolo sono proporzionali a tali lati. % lang-it
Ma una proprietà molto più importante è la seguente: % lang-it

\begin{proposition}
    \label{random_prp_symmedians_proportional_squares}
    Symediana poprowadzona z wierzchołka $C$ trójkąta $ABC$ dzieli wewnętrznie przeciwległy bok proporcjonalnie do kwadratów długości boków przyległych: % lang-pl
    La symmediana condotta dal vertice $C$ del triangolo $ABC$ divide internamente il lato opposto in proporzione ai quadrati delle lunghezze dei lati adiacenti: % lang-it
    \begin{equation}
        \frac{AF}{BF} = \frac{b^2}{a^2}.
    \end{equation}
\end{proposition}

I odwrotnie, punkt leżący na boku $AB$ i spełniający powyższą proporcję wyznacza symedianę z~wierzchołka $C$. % lang-pl
Zetel \cite[s. 111, 112]{zetel_2020} uważa to za szczególny przypadek twierdzenie Steinera o~prostych izogonalnych. % lang-pl
(Chyba nie ma tego twierdzenia w tej książce). % lang-pl
Viceversa, un punto appartenente al lato $AB$ che soddisfa la proporzione precedente determina la symmediana uscente dal vertice $C$. % lang-it
Zetel \cite[p. 111, 112]{zetel_2020} considera questo un caso particolare del teorema di Steiner sulle rette izogonali. % lang-it
(Probabilmente questo teorema non compare nel presente libro). % lang-it

\begin{proposition}
    Symediana wychodząca z wierzchołka $A$ trójkąta $ABC$ ma długość % lang-pl
    La symmediana uscente dal vertice $A$ del triangolo $ABC$ ha lunghezza % lang-it
    \begin{equation}
        \frac{bc}{b^2 + c^2} \cdot \sqrt {2b^2 + 2c^2 - a^2}.
    \end{equation}
\end{proposition}

Zetel \cite[s. 119]{zetel_2020} pisze, że można dostać tę długość z twierdzenia Stewarta. % lang-pl
Zetel \cite[p. 119]{zetel_2020} scrive che questa lunghezza può essere ottenuta mediante il teorema di Stewart. % lang-it
% ale nie trzeba.

\begin{proposition}
[punkt Lemoine'a] % lang-pl
[punto di Lemoine] % lang-it
    \label{punkt_lemoine}%
    W ustalonym trójkącie $ABC$, trzy symediany przecinają się w jednym punkcie $L$, zwanym punktem Lemoine'a. % lang-pl
    In un triangolo fissato $ABC$, le tre symmediane si incontrano in un unico punto $L$, detto punto di Lemoine. % lang-it
\end{proposition}

Zetel \cite[s. 114]{zetel_2020} sugeruje udowodnić \ref{punkt_lemoine} na podstawie twierdzenia odwrotnego do twierdzenia Cevy. % lang-pl
Ross Honsberger nazwie istnienie punktu Lemoine'a ,,jednym z klejnotów koronnych współczesnej geometrii''. % lang-pl
Zetel \cite[p. 114]{zetel_2020} suggerisce di dimostrare \ref{punkt_lemoine} mediante il teorema inverso di Ceva. % lang-it
Ross Honsberger definirà l'esistenza del punto di Lemoine come \emph{uno dei gioielli della corona della geometria moderna}. % lang-it

\begin{proposition}
    Punkt Lemoine'a minimalizuje sumę kwadratów odległości od boków trójkąta. % lang-pl
    Il punto di Lemoine minimizza la somma dei quadrati delle distanze dai lati del triangolo. % lang-it
\end{proposition}

% TODO: https://www.deltami.edu.pl/media/articles/2015/02/delta-2015-02-krotka-opowiesc-o-symedianie.pdf

\begin{proposition}
    Punkt Lemoine'a $L$ trójkąta $ABC$ jest środkiem ciężkości trójkąta, którego wierzchołki to rzuty punktu $L$ na boki $AB$, $BC$, $AC$. % lang-pl
    Il punto di Lemoine $L$ del triangolo $ABC$ è il baricentro del triangolo i cui vertici sono le proiezioni di $L$ sui lati $AB$, $BC$, $AC$. % lang-it
\end{proposition}

\begin{proposition}
[twierdzenie Schlömilcha] % lang-pl
[teorema di Schlömilch] % lang-it
    Trzy proste łączące środki boków trójkąta ze środkami odpowiednich wysokości są współpękowe, przechodzą przez punkt Lemoine'a. % lang-pl
    Le tre rette che congiungono i punti medi dei lati di un triangolo con i punti medi delle rispettive altezze sono concorrenti e passano per il punto di Lemoine. % lang-it
\end{proposition}

\index{twierdzenie!Schlömilcha} % lang-pl
\index{teorema!di Schlömilch} % lang-it

Napisze o tym duet Bogdańska, Neugebauer \cite[s. 195]{neugebauer_2018} (jako ćwiczenie, bez wzmianki o nazwie punktu), ale też Zetel \cite[s. 34]{zetel_2020}. % lang-pl
Ne scrivono Bogdańska e Neugebauer \cite[p. 195]{neugebauer_2018} (come esercizio, senza menzionare il nome del punto), ma anche Zetel \cite[p. 34]{zetel_2020}. % lang-it
% https://ems.press/content/serial-article-files/51635 VVV
Twierdzenie napotka w 1862 roku Oskar Schlömilch, cokolwiek to znaczy. % lang-pl
Il teorema fu scoperto da Oskar Schlömilch nel 1862, qualunque cosa ciò voglia dire. % lang-it

\begin{definition}
[antyrównoległa] % lang-pl
[antiparallela] % lang-it
    Antyrównoległa do boku $AB$ w trójkącie $ABC$ to prosta, która tnie bok naprzeciw wierzchołka $A$ pod takim samym kątem jak ten przy wierzchołku $A$ oraz bok naprzeciw wierzchołka $B$ pod takim samym kątem jak ten przy wierzchołku $B$. % lang-pl
    Un'antiparallela al lato $AB$ nel triangolo $ABC$ è una retta che interseca il lato opposto al vertice $A$ con lo stesso angolo presente in $A$ e il lato opposto al vertice $B$ con lo stesso angolo presente in $B$. % lang-it
\end{definition}

\begin{proposition}
    Trzy antyrównoległe do boków trójkąta przechodzące przez punkt Lemoine'a wyznaczają sześć nowych punktów na brzegu trójkąta. Leżą one na pierwszym okręgu Lemoine'a. % lang-pl
    Le tre antiparallele ai lati del triangolo passanti per il punto di Lemoine determinano sei nuovi punti sul contorno del triangolo. Essi giacciono sul primo cerchio di Lemoine. % lang-it
\end{proposition}

\begin{proposition}
    Trzy równoległe do boków trójkąta przechodzące przez punkt Lemoine'a wyznaczają sześć nowych punktów na brzegu trójkąta. Leżą one na drugim okręgu Lemoine'a. % lang-pl
    Le tre rette parallele ai lati del triangolo passanti per il punto di Lemoine determinano sei nuovi punti sul contorno del triangolo. Essi giacciono sul secondo cerchio di Lemoine. % lang-it
\end{proposition}

Istnieje jeszcze trzeci okrąg Lemoine'a, ale jego konstrukcja nie jest już taka prosta. % lang-pl
Esiste anche un terzo cerchio di Lemoine, ma la sua costruzione non è altrettanto semplice. % lang-it

%

\todofoot{Czewiany, w tym izotomiczne i izogonalne.}

\todofoot{Twierdzenie Steinera}

%